\documentclass[journal]{IEEEtran}
\usepackage{amsmath,amssymb}
\usepackage{booktabs}
\usepackage{graphicx}
\usepackage{hyperref}
\usepackage{url}
\title{Optimization Algorithm Atlas: A Mechanism--Evidence--Gap Framework for Reproducible Optimization Research}
\author{Mohammad Amir Khusru Akhtar\\
	\small Usha Martin University, Ranchi -- 834001, India\\
	\small \texttt{akakhtar.2024@gmail.com}}
\begin{document}
\maketitle
\begin{abstract}
Optimization research contains a large and rapidly evolving collection of methods whose names, mechanisms, theoretical guarantees, benchmark evidence, and reproducibility are often discussed at different levels. This paper develops an Optimization Algorithm Atlas that separates these layers and reconnects them through explicit provenance. A 90-entry cross-paradigm census is normalized into a representative mechanism layer containing 42 source-backed fingerprints linked to 46 verified source records. Evidence is indexed by algorithm and condition and classified as documented, partial, contradictory, untested, or project-validated, preventing absence of evidence from being interpreted as failure. The framework is coupled to an executable continuous black-box characterization using Random Search, differential evolution, Nelder--Mead, and CMA-ES on all 24 noiseless COCO BBOB functions across dimensions 5, 10, 20, and 40 under normalized evaluation budgets. Discovery and held-out instances, independent objective-call accounting, and COCO-authoritative target logging provide an auditable path from protocol to claim. The held-out campaign supports condition-level characterization but not a post hoc binary failure frontier because its exact reliability threshold was not frozen before confirmation. Consequently, no new optimizer is introduced. The resulting atlas provides a mechanism-first framework for evidence synthesis, comparator selection, benchmark design, novelty checking, reproducibility auditing, and explicit research evidence obligations.
\end{abstract}

\begin{IEEEkeywords}
optimization algorithms, algorithm taxonomy, benchmarking, reproducibility, evidence synthesis, black-box optimization
\end{IEEEkeywords}
\section{Introduction}
\label{sec:introduction}

Optimization research spans mathematical programming, gradient and derivative-free methods, evolutionary and swarm search, distribution-based optimization, surrogate methods, multiobjective optimization, robust and stochastic formulations, dynamic and online optimization, and learning-assisted approaches. At the same time, the number of named metaheuristics has grown rapidly. Large reviews have catalogued hundreds of methods, while recent structural studies have shown that differently named algorithms can share substantial computational structure. Functional taxonomies and benchmark-driven reviews have further shifted attention from inspiration metaphors toward search dynamics and empirical behavior.

These developments expose a second problem beyond algorithm proliferation: evidence is fragmented across incompatible levels. A canonical paper may define an update rule, a theoretical study may establish convergence under specific assumptions, a benchmark paper may report finite-budget behavior, a later variant may introduce adaptation, and an implementation may differ from the published specification. Broad statements about scalability or robustness can therefore mix theory, implementation, benchmark distribution, parameter policy, and computational budget unless their conditions are retained.

This paper addresses that evidence problem through a mechanism--evidence--gap atlas. Algorithm identity, computational mechanism, empirical evidence, reproducibility, and unresolved questions are represented separately before synthesis. An algorithm name is retained for provenance; a mechanism fingerprint describes search state, information, proposal geometry, memory, selection, adaptation, diversity, constraints, and models; an evidence state records what is known under a declared condition; and an evidence obligation states what observation would resolve an uncertainty.

The framework is cross-paradigm. The census contains 90 optimization identities spanning classical, derivative-free, evolutionary, swarm, probabilistic, surrogate, multiobjective, constrained, robust, dynamic, mixed-variable, and learning-assisted optimization. These identities are not interpreted as 90 distinct mechanisms. A representative evidence layer contains 42 source-backed mechanism fingerprints linked to 46 verified source records. This separation preserves broad literature coverage without equating naming diversity with mechanistic novelty.

The paper also distinguishes evidence absence from algorithmic failure. Condition-level states are represented as documented, partial, contradictory, untested, or project-validated. Structural similarity can generate a hypothesis but cannot transfer empirical evidence between algorithms. A missing experiment creates an evidence obligation rather than proof of a defect.

To connect synthesis with executable evidence, the atlas includes a frozen continuous black-box characterization using the noiseless COCO BBOB suite. Random Search, SciPy differential evolution, bounded SciPy Nelder--Mead, and pycma CMA-ES are evaluated under normalized objective-evaluation budgets across dimensions 5, 10, 20, and 40. Discovery and held-out instance sets are separated, objective-call accounting is cross-checked against COCO, and official \texttt{cocopp} postprocessing is retained for target-runtime analysis. The experiment demonstrates condition-level evidence generation rather than a universal optimizer ranking.

An important outcome is methodological restraint. The held-out campaign executed successfully, but the exact numerical reliability threshold and target subset required for a binary failure-frontier claim were not frozen before confirmation. The paper therefore retains the campaign as project-validated characterization and does not construct an exact frontier afterward. No new optimizer is introduced because the frozen evidence does not establish the independently supported mechanism-level repair opportunity required to justify one.

The contributions are fourfold. First, the paper develops a mechanism-first representation separating algorithm identity from state, information, proposal, memory, selection, adaptation, and problem-structure mechanisms. Second, it organizes heterogeneous literature into condition-level evidence states linked to provenance. Third, it couples the review to an executable reproducibility chain from protocol and configuration through frozen artifacts and claims. Fourth, it converts unresolved observations into explicit evidence obligations and a researcher decision framework for comparator selection, benchmark design, novelty checking, and future experiments.

The objective is not to replace specialized optimization theory or domain-specific benchmarking. It is to provide a common evidence architecture through which those results can be compared without erasing their assumptions and conditions. The remainder of the paper defines the review methodology, develops the mechanism taxonomy, constructs the evidence atlas, separates theory from finite-budget evidence, presents the controlled computational characterization, evaluates gaps and existing remedies, and concludes with a decision framework and bounded research frontier.

\section{Review Questions, Scope, and Evidence Methodology}
\label{sec:methods}

\subsection{Review objective and research questions}

This study treats optimization research as an evidence-organization problem rather than as a catalogue of algorithm names. The review asks which algorithm families have been studied, which computational mechanisms distinguish them, what theoretical and empirical evidence supports their use, how benchmark design conditions the conclusions that can be drawn, which implementations are reproducible, and where the available evidence remains incomplete. Particular attention is given to structural overlap among nominally different methods, condition-dependent performance deterioration, and the distinction between documented limitations and genuinely unresolved research gaps.

Nine review questions organize the analysis. They concern algorithm and problem-class coverage; mechanism-level distinctions; evidence on quality, convergence, robustness, scalability, computational cost, and reliability; benchmark adequacy; structural similarity; condition-dependent failure regimes; reproducibility; evidence-supported gaps; and the circumstances under which a gap would justify an established remedy, a transferable mechanism, or a new optimizer. A new optimizer is therefore not a predetermined outcome of the review.

\subsection{Evidence classes and claim discipline}

Four evidence classes are kept separate throughout the study: documented facts supported directly by a source; empirical results reported by prior studies; cross-study synthesis produced by the review; and project-generated computational findings. Literature-reported evidence is not relabelled as project validation, and project results are not used to strengthen claims beyond the conditions actually evaluated.

The review uses a 90-entry algorithm census as an identity and coverage layer. The census is not interpreted as evidence for 90 distinct mechanisms. Algorithms are instead normalized into a mechanism-first representation. Following the final evidence audit, the representative mechanism layer contains 42 source-backed fingerprints linked to 46 verified source records. Named variants are retained when they introduce a relevant structural distinction; nominally different methods are not promoted to separate mechanisms merely because they use different metaphors or terminology.

\subsection{Search, screening, and evidence extraction}

The evidence-mining protocol covers exact and mathematical optimization, gradient and proximal methods, derivative-free search, evolutionary and population-based optimization, swarm methods, probabilistic and distribution-based search, Bayesian and surrogate optimization, multi- and many-objective optimization, constrained and robust optimization, dynamic and online optimization, mixed-variable and combinatorial methods, bilevel optimization, and learning-assisted optimization.

Search concepts combine optimization terminology with mechanism, benchmark, scalability, robustness, constraint, noise, dynamics, mixed-variable, reproducibility, structural-bias, comparison, sensitivity, ablation, and parameter-control terms. Sources are eligible when they contribute an original algorithm or substantial mechanism, a theoretical result, benchmark methodology, a systematic or critical synthesis, comparative evidence, reproducibility evidence, structural-bias analysis, benchmark-suite design, negative evidence, or a real-world evaluation. Sources that do not provide an inspectable algorithmic definition, cannot be linked to a verifiable bibliographic record, report selected best runs without adequate context, or use comparisons whose budgets cannot be interpreted fairly are excluded or isolated rather than silently incorporated.

For each representative algorithm, evidence extraction records identity and provenance, representation, initialization, information sources, update geometry, selection or replacement, memory, adaptation, diversity or restart mechanisms, constraint handling, derivative information, surrogate use, stopping conditions, theoretical guarantees and assumptions, benchmark coverage, parameter policy, statistical practice, implementation availability, and known limitations. Untested conditions are recorded as evidence gaps; absence of evidence is not interpreted as algorithmic failure.

\subsection{Mechanism-first normalization}

Each representative method is encoded using a common structural fingerprint comprising state type, information source, proposal geometry, memory, adaptation, selection, diversity control, constraint handling, variable scope, objective scope, derivative order, model type, parallelism, dominant overhead, and a documented or hypothesized failure signal. This normalization provides a common language for comparing methods that originate in different research traditions.

Structural similarity is interpreted as similarity of encoded computational mechanisms, not as algorithmic identity, plagiarism, or equivalent empirical performance. Similarity analyses are therefore used as an exploratory synthesis tool and are checked against known genealogical relationships and manual mechanism inspection.

\subsection{Benchmark and comparison protocol}

The computational component uses the noiseless COCO BBOB suite and retains all 24 functions rather than selecting functions after observing results. The core design considers dimensions \(D\in\{2,5,10,20,40\}\). The principal cross-dimensional campaigns use dimensions 5, 10, 20, and 40. Discovery uses BBOB instances 1--5, while a separately frozen held-out characterization uses instances 6--10. Stochastic configurations use seeds 11, 23, and 37. Evaluation budgets are normalized as \(100D\), \(300D\), and \(1000D\) objective evaluations.

The maintained comparator panel comprises Random Search as a transparent reference baseline, SciPy differential evolution with an explicitly frozen configuration and polishing disabled, bounded SciPy Nelder--Mead as a direct-search comparator, and pycma CMA-ES through an ask--tell interface. Algorithms whose publication implementation or configuration provenance did not pass the same acceptance process were not inserted merely to enlarge the comparison panel.

All publication computations use objective-level accounting. The objective recorder increments once for each actual objective call, maintains best-so-far values, records declared checkpoints, and is cross-checked against COCO's evaluation counter. Evaluation budgets are therefore defined in objective calls rather than generations or optimizer-specific iterations. Failures and unobserved checkpoints are retained rather than imputed as successful observations.

\subsection{Statistical analysis and interpretation}

The atomic experimental unit is algorithm/configuration \(\times\) benchmark function \(\times\) instance \(\times\) dimension \(\times\) seed \(\times\) budget. Primary views include target attainment, evaluations to target, best-so-far values at declared budgets, anytime or empirical cumulative distribution views where supported by suite semantics, and computational overhead as a separate cost dimension. Aggregate ranks, when reported, are descriptive and are not used to establish universal superiority.

Comparisons are matched on function, instance, dimension, seed, and budget whenever the algorithms permit such pairing. Distributional summaries and confidence intervals accompany the comparisons. Paired effect estimates with bootstrap confidence intervals are used where appropriate, and multiplicity-controlled post-hoc comparisons follow an omnibus analysis when inferential comparisons are warranted. Holm-type family-wise error control is the default for planned post-hoc families unless the endpoint structure requires a separately declared procedure.

\subsection{Discovery, held-out characterization, and frontier claims}

Discovery and confirmation data are separated to prevent a condition identified from the data from serving as its own sole confirmation. The EF003 discovery campaign uses instances 1--5. EF005 repeats the frozen algorithms, dimensions, seeds, and evaluation budgets on held-out instances 6--10. The computational and provenance pipelines, including algorithm-separated COCO logging and joint \texttt{cocopp} postprocessing, passed the frozen validation workflow.

The original frontier protocol required an exact target subset and a numerical reliability threshold to be fixed before held-out outcomes were inspected. Those exact numerical choices were not frozen before the held-out campaign was executed. Consequently, the paper does not promote a threshold-dependent binary failure frontier after the fact. The held-out campaign is used for condition-level target-runtime and scaling characterization and for assessing whether qualitative patterns persist on disjoint instances. No post hoc exact boundary, universal failure claim, or mechanism-repair claim is inferred solely from performance deterioration.

\subsection{Gap validation and remedy rule}

Evidence gaps are classified as theoretical, mechanistic, structural-bias, benchmark, scalability, robustness, constraint, dynamic, mixed-variable, reproducibility, comparison, real-world, statistical, or computational-cost gaps. A gap record distinguishes what is documented, partially supported, contradictory, untested, or project-validated.

A new mechanism or optimizer is considered only when evidence supports a mechanism-level repair opportunity that cannot be adequately explained by tuning, budget normalization, an existing adaptive variant, restart, established constraint handling, or another known remedy. Because the frozen evidence does not establish such a confirmatory G4/G5 repair opportunity, no new optimizer or result-driven ablation campaign is introduced in this paper.

\subsection{Reproducibility and evidence freeze}

The research workflow separates calibration, pilot, rehearsal, run-complete-unverified, and project-validated states. Calibration, pilot, and rehearsal outputs are never cited as publication evidence. Benchmark software, comparator configurations, target semantics, statistical rules, execution stacks, and interpretation boundaries are versioned in the repository.

Before manuscript drafting, the P031 evidence-freeze workflow validated the execution-freeze records, literature/source integrity, final pre-freeze audit, and a Git-blob manifest of the principal evidence and publication-control artifacts. The manuscript is written against this frozen boundary. Subsequent scientific changes require a separately versioned amendment and cannot silently replace the evidence supporting the reported conclusions.

\section{Mechanism Taxonomy and Genealogy}
\label{sec:taxonomy}

\subsection{From algorithm names to computational mechanisms}

Optimization taxonomies are often organized by historical labels, mathematical assumptions, or the metaphor used to introduce an algorithm. These views remain useful for locating a method in the literature, but they can obscure structural relationships. Two algorithms with different names may preserve the same kind of search state, use similar information, and generate proposals through closely related operators; conversely, algorithms grouped under one broad family may differ substantially in memory, adaptation, selection, or model structure. We therefore use the algorithm census as an identity layer and a separate mechanism fingerprint as the unit of structural comparison.

The taxonomy is organized around three questions: what state is carried between evaluations, what information is used to update that state, and how is a new candidate or decision generated? Secondary axes describe memory, adaptation, selection, diversity control, constraint handling, derivative order, model type, parallelism, and dominant computational overhead. This representation does not assert that these fields completely determine algorithm behavior. Rather, it provides a common descriptive basis on which algorithms from otherwise separate literatures can be compared.

\subsection{State, information, and proposal geometry}

The most immediate distinction concerns the state maintained by the optimizer. Point-based methods maintain one incumbent or iterate, possibly augmented by a limited history. Nelder--Mead evolves a simplex, while DIRECT maintains a geometric partition of the domain. Population methods retain multiple candidate solutions and use interactions among them. Distribution-based methods, including CMA-ES, natural evolution strategies, estimation-of-distribution algorithms, and the cross-entropy method, elevate a parameterized sampling distribution to a primary search state. Surrogate methods maintain observations together with an explicit predictive or density model. ADMM maintains coupled primal and dual state, while sample-average approximation retains a sampled scenario representation of an uncertain objective.

A second axis describes the information transformed into search moves. Gradient and proximal methods use local differential information, with proximal operators incorporating composite structure. L-BFGS augments gradients with a short curvature history, coordinate descent restricts updates to selected coordinates, and online gradient descent uses the gradient of the currently revealed loss. Derivative-free local search instead obtains direction from objective comparisons, simplex geometry, polling directions, interpolation models, or neighborhood structure.

Population methods construct direction from relations among evaluated candidates. Differential evolution uses pairwise population differences; PSO combines personal and social best information through velocity dynamics; genetic algorithms use selected parents followed by recombination and mutation; ant-colony methods use accumulated pheromone information to bias constructive decisions. These mechanisms should not be collapsed into a single population-heuristic category because the information converted into a proposal is materially different.

Distribution-based search provides another transition. CMA-ES updates a covariance-structured search distribution using ranked samples and evolution paths. Natural evolution strategies update distribution parameters through a natural-gradient view, whereas estimation-of-distribution algorithms and the cross-entropy method refit probabilistic models from selected or elite samples. Surrogate optimization replaces or supplements direct search geometry with a model of observed performance. Gaussian-process Bayesian optimization uses posterior-dependent acquisition, TPE uses conditional density models, SMAC uses a model suited to mixed configuration spaces, and TuRBO localizes Bayesian optimization through trust regions. Hyperband is structurally different because its principal mechanism is resource allocation and successive elimination; BOHB combines resource allocation with model-based sampling.

\subsection{Selection, memory, and adaptation}

Selection determines which observations influence future search. Greedy replacement in differential evolution differs from probabilistic acceptance in simulated annealing, Pareto and diversity selection in multiobjective evolutionary algorithms, indicator selection in SMS-EMOA, and acquisition-driven selection in Bayesian optimization. Feasibility rules alter the ordering itself by incorporating constraint violation.

Memory ranges from an incumbent or short gradient history to archives, personal-best records, tabu lists, pheromone matrices, parameter memories, and learned probabilistic models. Tabu memory suppresses recently visited moves and supports diversification; personal-best memory in PSO influences velocity; the external archive in JADE changes the pool from which differential information can be drawn; SHADE stores successful control-parameter information; Bayesian optimization stores observations to update a surrogate.

Adaptation is therefore described by what is adapted and from which evidence. The DE lineage makes this distinction visible. Canonical DE uses externally selected control parameters. jDE self-adapts control parameters with individuals. SaDE uses accumulated search experience to adapt strategy choice and parameters. JADE introduces success-based parameter adaptation together with a current-to-p-best strategy and optional archive. SHADE adds an explicit history of successful parameter settings, and L-SHADE further introduces scheduled population-size reduction. These methods share differential variation while modifying the memory and control loops surrounding it.

\subsection{Constraints, uncertainty, and temporal structure}

Constraint handling is not a single auxiliary operation. COBYLA constructs local linear approximations of the objective and constraints inside a trust-region process. Feasibility rules alter dominance between candidates without requiring a problem-specific penalty coefficient. ADMM exposes constraints through operator splitting and coupled primal--dual updates. Robust optimization embeds uncertainty sets into the decision model, whereas sample-average approximation replaces an expectation by a finite scenario average. These mechanisms address different information structures and should not be compared solely under a generic constrained-optimization label.

Temporal information creates another distinction. Online gradient descent receives losses sequentially and is naturally evaluated through regret rather than only terminal objective value. Dynamic optimization methods face an objective or environment that changes over time. The census retains such research families even when the representative evidence layer does not promote every named dynamic method to a distinct fingerprint. Census coverage and mechanism evidence are therefore deliberately separate.

\subsection{Multiobjective selection as a structural layer}

Multiobjective optimization demonstrates that proposal generation alone is insufficient to characterize an algorithm. NSGA-II uses nondominated sorting and crowding to combine convergence and diversity. SPEA2 uses strength-based fitness, density estimation, and archive truncation. MOEA/D converts a multiobjective problem into linked scalar subproblems. NSGA-III introduces reference-point-based diversity management for many-objective settings, RVEA uses reference vectors, and SMS-EMOA uses hypervolume contribution as a selection signal. These algorithms may share conventional variation operators while differing principally in environmental selection and diversity geometry.

Accordingly, objective scope, selection, memory, and diversity control are encoded separately. This avoids treating all Pareto-based evolutionary methods as structurally interchangeable and permits mechanisms such as decomposition, reference directions, and indicators to be compared independently of the variation operator with which they are implemented.

\subsection{Genealogy versus structural similarity}

Genealogy and structural similarity answer different questions. Genealogy records documented historical or methodological descent: jDE, SaDE, JADE, SHADE, and L-SHADE belong to the broader DE development line, while BOHB explicitly combines model-based optimization with Hyperband-style resource allocation. Structural similarity instead asks whether normalized mechanism fields are close, regardless of documented descent. High structural similarity therefore does not establish ancestry, equivalence, or equivalent empirical behavior.

This distinction is particularly important for metaphor-driven optimization. The census retains a redundancy-audit lane containing methods introduced through different natural or social metaphors, but the evidence framework does not assume that each name defines a distinct mechanism. Promotion to a separate mechanism fingerprint requires an inspectable structural distinction supported by evidence. This rule preserves literature coverage without equating naming proliferation with mechanistic novelty.

\subsection{A mechanism-loop representation}

The fingerprints can be summarized through the descriptive update
\begin{equation}
\mathcal{S}_{t+1}=\mathcal{M}\!\left(\mathcal{S}_{t},\mathcal{I}_{t},\mathcal{P}(\mathcal{S}_{t},\mathcal{I}_{t};\theta_t),\mathcal{E}_{t}\right),
\label{eq:mechanism-loop}
\end{equation}
where \(\mathcal{S}_{t}\) is persistent search state, \(\mathcal{I}_{t}\) is information available at decision time \(t\), \(\mathcal{P}\) is the proposal mechanism, \(\theta_t\) is its control state, \(\mathcal{E}_{t}\) is newly observed objective, constraint, gradient, or model evidence, and \(\mathcal{M}\) is the memory, selection, and adaptation update. This notation is descriptive rather than a convergence theorem; it exposes which component changes when named algorithms are compared.

Under this view, algorithm development can change the state representation, information source, proposal geometry, evidence-retention rule, selection criterion, adaptation policy, or constraint and uncertainty model. A claimed new optimizer can therefore be examined at the component level. The same representation links the taxonomy to the evidence analysis that follows: empirical limitations can be associated with mechanism components without assuming that finite benchmark evidence establishes an intrinsic or universal failure.

\section{Condition-Level Evidence Atlas}
\label{sec:evidence-atlas}

\subsection{Evidence is indexed by condition, not by winner}

The atlas is designed to answer a different question from a benchmark leaderboard. For an algorithm \(A\), an empirical statement is indexed by a condition vector \(\theta\) describing the problem, dimension, budget, instance distribution, parameter policy, implementation, and evaluation criterion. Evidence is represented conceptually as
\begin{equation}
\mathcal{E}(A,\theta)=\{s,c,p\},
\label{eq:evidence-state}
\end{equation}
where \(s\) is the evidence state, \(c\) its evidential strength, and \(p\) the provenance linking the statement to a source or frozen project artifact. A result obtained under one benchmark regime is therefore not silently generalized to another.

Five states are distinguished: documented, partial, contradictory, untested, and project-validated. Documented evidence directly supports the stated condition. Partial evidence supports only part of the condition or claim. Contradictory evidence records materially different findings requiring further resolution. Untested denotes absence of adequate evidence and is not synonymous with failure. Project-validated is reserved for repository-generated findings that passed the frozen provenance and validation workflow.

\subsection{Canonical provenance and condition-level evidence}

The representative mechanism layer contains 42 source-backed fingerprints linked to 46 verified source records. Provenance for defining mechanisms is therefore denser than evidence about the conditions under which particular limitations appear. A well-established algorithmic definition does not imply that its scaling, robustness, failure transitions, parameter sensitivity, or computational overhead have been comprehensively mapped.

The condition-level status matrix illustrates this distinction. High-dimensional simplex dynamics for Nelder--Mead and nonconvex curvature behavior for BFGS are represented as partial evidence rather than universal failure claims. The nonsmooth-objective limitation associated with L-BFGS is documented within known classes, while extension beyond those classes remains a mapping problem. Simulated annealing retains a partial finite-budget cooling question, and premature pheromone concentration in ant-colony optimization remains a partial mechanism-level concern requiring controlled study.

The multiobjective records follow the same discipline. Many-objective crowding behavior for NSGA-II is documented at the mechanism level, while a controlled diagnostic geometry map remains useful. Evidence concerning difficult Pareto-front geometries for MOEA/D and NSGA-III remains partial. Each status therefore applies to a mechanism--condition pair, not to an algorithm as a whole.

\subsection{Absence of evidence is explicit}

The census covers a broader space than the condition-level status matrix. Dynamic, mixed-variable, bilevel, learning-assisted, robust, stochastic, constrained, and combinatorial optimization remain unevenly characterized. Many algorithm--condition combinations have canonical definitions but no project-validated failure map. This asymmetry is retained rather than filled by inference.

The distinction between untested and failed is fundamental. If controlled evidence is unavailable for an algorithm under a condition, the atlas records an evidence obligation rather than a negative performance judgment. A missing study can establish an evidence gap; it cannot by itself establish a mechanistic defect.

\subsection{Benchmark coverage is part of the evidence statement}

Benchmark suites are treated as measurement instruments with coverage limits. The publication experiment uses noiseless COCO BBOB as a controlled continuous black-box domain. Function identity, dimension, instance, and evaluation budget are retained rather than pooled away. Property-conditioned synthesis is permitted only when an authoritative mapping supports the grouping and within-group evidence is sufficient.

Noisy and discrete domains remain separate. COCO BBOB-noisy and IOH pseudo-Boolean optimization appear in the suite registry as distinct evidence domains, but registry or adapter availability does not promote them into the present project-validated publication evidence. Local Sphere, Rastrigin, and Rosenbrock problems are calibration instruments and are excluded from publication evidence unless separately promoted. Dimension and budget are scientific axes, not nuisance metadata; a method is therefore not summarized by one unconditional rank across \(D\) and FE/D conditions.

\subsection{Reproducibility as an evidence dimension}

Reproducibility is not inferred from an algorithm name or canonical paper. The reproducibility registry separates definition availability, pseudocode or specification, reference implementation, independent implementation, parameter disclosure, random-seed disclosure, benchmark version, environment disclosure, and rerun status. Fields that remain pending are not converted into positive reproducibility scores.

The project-generated BBOB evidence follows an executable provenance path. Software versions, configurations, suite, instances, seeds, budgets, target semantics, objective-call accounting, and interpretation rules were frozen before promotion. COCO logging provides target-runtime authority, while an independent objective recorder provides exact objective-call and checkpoint accounting. The counters were cross-checked before publication use. Reproducibility is thus treated as a property of an evidence-producing workflow rather than a reputation attached to an algorithm.

\subsection{Structural similarity does not transfer evidence}

Mechanism fingerprints identify structural proximity, but empirical evidence is not automatically transferred between neighboring algorithms. The DE lineage shares differential variation, yet jDE, SaDE, JADE, SHADE, and L-SHADE modify parameter adaptation, strategy control, archives, success-history memory, or population-size schedules. Evidence of a limitation in one member can motivate a hypothesis about another, but cannot establish the same failure regime.

The restriction also applies across broader classes. Distribution-based search links CMA-ES, natural evolution strategies, EDAs, and the cross-entropy method, but their update geometry and assumptions differ. Surrogate methods share a model layer, yet Gaussian-process Bayesian optimization, TPE, SMAC, and TuRBO encode different response assumptions and search locality. Structural similarity therefore generates comparisons and evidence obligations; it does not manufacture missing empirical results.

\subsection{Negative and contradictory evidence}

Negative evidence is retained without converting it into a universal judgment. Reports of premature convergence or exploration--exploitation limitations in success-history DE variants, for example, motivate condition-specific questions about diversity and adaptation but do not establish a defect across the entire DE lineage.

Contradictory findings are treated as unresolved conditional evidence. Differences in benchmark distribution, implementation, tuning, evaluation budget, stopping rule, or statistical unit can produce disagreement without implying that one study is invalid. The atlas records the unresolved condition and the minimal experiment needed to distinguish explanations. Contradiction thereby becomes a research object rather than a reason to select the result supporting a preferred narrative.

\subsection{From gaps to resolving experiments}

An evidence gap is promoted only when it can be stated operationally: what is known, what is missing, why the missing evidence matters, the strongest provenance, a testable hypothesis, and a minimal resolving experiment. The resulting discipline is
\begin{equation}
\text{missing evidence}\;\not\Rightarrow\;\text{mechanism failure}\;\not\Rightarrow\;\text{new optimizer}.
\label{eq:gap-discipline}
\end{equation}
Missing evidence first creates an evidence obligation. A controlled experiment may establish a condition-level limitation. Only a reproduced limitation with an independently supported mechanism explanation can motivate a repair hypothesis. The frozen project did not reach that final repair threshold, so no result-driven optimizer is introduced.

\subsection{Project-generated characterization}

The project experiments complement the literature atlas with a standardized continuous black-box characterization under one frozen protocol. They do not replace the heterogeneous literature and do not establish a universal ordering. Random Search, SciPy differential evolution, bounded SciPy Nelder--Mead, and pycma CMA-ES are evaluated on matched BBOB conditions under normalized objective-evaluation budgets. Discovery instances 1--5 and held-out instances 6--10 are separated.

Because the exact numerical reliability threshold and target subset required for a confirmatory binary failure frontier were not frozen before the held-out campaign, the project evidence is used as condition-level characterization rather than as a post hoc exact boundary. The strength of a claim is determined not only by observed data, but also by whether the decision rule supporting that claim was specified before the relevant evidence was examined.

\section{Theory, Guarantees, Failure Modes, and Reproducibility}
\label{sec:theory-repro}

\subsection{Guarantees are conditional statements}

A theoretical guarantee is meaningful only together with its assumptions. Smoothness, convexity, Lipschitz conditions, boundedness, stochastic assumptions, constraint qualifications, sampling assumptions, and access to derivatives or exact subproblem solutions determine what a theorem establishes. The atlas therefore does not reduce theory to a binary has-guarantee field. Instead, a theoretical statement is interpreted as a tuple
\begin{equation}
\mathcal{T}=\left(\mathcal{A},\mathcal{G},\mathcal{O}\right),
\end{equation}
where \(\mathcal{A}\) denotes assumptions, \(\mathcal{G}\) the guaranteed property, and \(\mathcal{O}\) the object to which the guarantee applies, such as iterates, expected regret, stationarity, convergence in probability, or approximation quality.

This distinction prevents guarantees from being compared outside their domains. A convergence result for a smooth convex objective, a global-convergence statement for a derivative-free polling framework, an asymptotic stochastic-search result, and a regret bound for online optimization answer different questions. None is converted into an unconditional empirical superiority claim.

\subsection{Theory and finite-budget behavior are complementary}

Asymptotic guarantees do not determine finite-budget performance, and strong finite-budget benchmark results do not establish a theorem. The evidence model therefore keeps theoretical and empirical axes separate. An algorithm may satisfy a convergence result while exhibiting unfavorable scaling under a particular budget, or may perform strongly on a benchmark despite having weaker available theory.

The distinction is especially important for stochastic and population-based optimization. Infinite-time convergence arguments, where available, need not predict behavior at \(100D\), \(300D\), or \(1000D\) evaluations. Conversely, a finite BBOB comparison cannot establish global convergence or refute a theorem whose assumptions and limiting regime differ from the experiment.

\subsection{Failure modes are mechanism--condition relations}

A failure mode is not defined as an algorithm being bad. It is a reproducible or documented relation between a mechanism and a condition under which an intended search behavior becomes ineffective, unstable, excessively costly, or insufficiently informative. We write this conceptually as
\begin{equation}
\mathcal{F}=\left(M,\theta,Y,P\right),
\end{equation}
where \(M\) is the implicated mechanism, \(\theta\) the condition, \(Y\) the observed outcome, and \(P\) its provenance. A mechanistic explanation requires more evidence than an observed deterioration.

The frozen literature matrix consequently uses bounded descriptions. High-dimensional simplex dynamics for Nelder--Mead, nonconvex curvature for BFGS, finite-budget cooling for simulated annealing, pheromone concentration for ant-colony optimization, and difficult Pareto-front geometries for decomposition or reference-direction methods are represented only at their supported evidence levels. Documented limitations such as known nonsmooth L-BFGS behavior are not generalized beyond the classes for which the evidence applies.

\subsection{From observation to mechanism claim}

The manuscript follows a claim-strength ladder. Descriptive algorithm identity occupies the lowest level and requires a canonical specification. Source-reported observations require the source and its conditions. Review synthesis requires comparable evidence across sources. A documented boundary requires a formal counterexample or adequate replicated empirical evidence. Project-validated conditional results require reproducible raw results, statistical analysis, and artifact provenance. A general mechanism claim requires cross-family evidence, tests of competing explanations, ablations, and external validation.

This ordering is important because performance deterioration alone cannot identify its cause. A higher-dimensional loss of target attainment could arise from search geometry, parameter scaling, population size, implementation overhead, termination, or the benchmark distribution. A mechanism explanation becomes credible only when alternatives are controlled or independently excluded. The present paper therefore does not promote its BBOB characterization to a general cross-family mechanism theory.

\subsection{Target semantics and authoritative measurement}

For the BBOB characterization, target attainment follows the frozen definition
\begin{equation}
f(x)\leq f_{\mathrm{opt}}+\Delta f,
\end{equation}
with the reporting grid \(\Delta f\in\{10^{2},10^{1},\ldots,10^{-8}\}\). The primary cross-function runtime unit is function evaluations divided by dimension, FE/D. Unreached targets remain unsuccessful observations rather than being removed from denominators.

The project deliberately separates measurement authorities. The COCO observer and logger are authoritative for benchmark target attainment and COCO-compatible postprocessing. The project objective recorder is authoritative for optimizer-independent objective-call accounting, best-so-far checkpoints, runtime metadata, and failures. The pinned public COCO Python interface does not expose a public \(f_{\mathrm{opt}}\) field suitable for a private reconstruction of exact target runtimes. Consequently, the analysis does not depend on undocumented internal optimum attributes or a hand-copied optimum table.

This split protects two different invariants: official benchmark semantics for target-runtime analysis and exact objective-call accounting for fixed-budget traces. Agreement between the recorder and COCO evaluation counts was checked in the publication executions.

\subsection{Failure-frontier discipline}

A failure frontier is useful only if its defining event is specified independently of the observations used to confirm it. Conceptually, for a performance requirement \(R\), a conditional frontier can be written
\begin{equation}
\partial\{\theta:R(A,\theta)=1\},
\end{equation}
where \(\theta\) indexes conditions such as dimension and budget. This notation does not by itself establish that such a frontier has been identified.

The project separated discovery instances 1--5 from held-out instances 6--10 and froze algorithms, dimensions, seeds, and FE/D budgets. However, the exact numerical reliability threshold and target subset required to convert target attainment into a binary frontier event were not frozen before the held-out campaign. Selecting those values afterward would turn a confirmatory boundary into a post hoc construction. The paper therefore reports the held-out campaign as project-validated characterization but does not claim an exact binary failure frontier.

This decision is a reproducibility result in its own right: experimental execution can be valid while a stronger interpretation remains unsupported because a decision rule was incompletely predeclared.

\subsection{Executable reproducibility}

Reproducibility is organized as a chain from scientific specification to executable evidence:
\begin{equation}
\text{protocol}\rightarrow
\text{configuration}\rightarrow
\text{execution freeze}\rightarrow
\text{raw artifact}\rightarrow
\text{analysis}\rightarrow
\text{claim}.
\end{equation}
Each link must remain inspectable. The repository therefore freezes benchmark software, comparator configurations, dimensions, instances, seeds, budgets, target semantics, statistical rules, and execution identifiers.

The production sequence also preserves historical corrections. Earlier execution freezes are not silently rewritten when a metadata or attribution defect is found. Instead, corrected freezes are versioned separately. This makes it possible to distinguish a scientific configuration change from a provenance or metadata correction.

For R10, the corrected cross-dimensional analysis contains 1,152 dimension--algorithm--function--budget cells. The EF004 production campaign produced 14,400 checkpoint rows and 64 algorithm-labelled COCO result roots across dimensions 5, 10, 20, and 40; objective-recorder and COCO counts agreed, and joint official \texttt{cocopp} postprocessing completed. The held-out EF005 campaign repeats the frozen design on disjoint instances for characterization.

\subsection{Evidence promotion and immutable freeze}

Not every executable output is publication evidence. Calibration, pilot, rehearsal, and run-complete-unverified states remain distinct from project-validated evidence. Failed, timed-out, and unobserved outcomes are retained rather than silently discarded. Statistical tables and figures are required to derive from promoted raw artifacts rather than hand-edited summaries.

Before manuscript drafting, the R15 audit checked execution-freeze integrity, source/fingerprint consistency, required interpretation records, and publication-control files. P031 then froze the principal evidence boundary using Git blob identities and validated that manifest in CI. Manuscript claims are consequently constrained by an immutable evidence snapshot.

The freeze does not imply that the atlas is complete. It means that the evidence supporting this version of the paper is fixed and auditable. A future experiment, corrected interpretation, or expanded literature audit must be introduced as a separately versioned amendment rather than silently changing the evidence beneath an existing claim.

\subsection{What reproducibility does not establish}

A reproducible experiment can still have limited external validity. Exact rerun capability does not prove that the chosen algorithms represent all optimization paradigms, that BBOB represents every application domain, or that a frozen parameterization is optimal for each method. Similarly, a canonical implementation does not guarantee fair computational cost across algorithms.

The paper therefore treats reproducibility as necessary for strong empirical claims but not sufficient for generality. Generalization depends additionally on benchmark coverage, mechanism diversity, parameter policy, cost accounting, independent replication, and consistency across conditions. This distinction motivates the bounded interpretation of the computational results in the next section.

\section{Controlled Computational Characterization}
\label{sec:computational}

\subsection{Experimental question}

The computational study asks a narrower question than a conventional leaderboard: how does a mechanism-diverse set of optimizers behave across a controlled continuous black-box test space when problem identity, dimensionality, instance set, seed policy, and objective-evaluation budget are fixed in advance? The retained panel comprises Random Search, SciPy differential evolution (DE), bounded SciPy Nelder--Mead (NM), and pycma CMA-ES. These methods were chosen because their implementations and frozen configurations passed the publication acceptance path; the panel is intentionally smaller than the census.

All 24 noiseless BBOB functions are evaluated at dimensions \(D\in\{5,10,20,40\}\). Discovery uses instances 1--5 and held-out characterization uses instances 6--10. Stochastic configurations use seeds 11, 23, and 37, with checkpoints at \(100D\), \(300D\), and \(1000D\) objective evaluations. ObjectiveRecorder counts were cross-checked against COCO evaluation counts.

\subsection{Discovery and held-out COCO evidence}

Figure~\ref{fig:coco_discovery_heldout} places the official COCO ECDF views for discovery and held-out runs side by side for the 20-D all-function aggregate. The two panels are generated directly by \texttt{cocopp} from EF004 and EF005 artifacts, respectively; they are not reconstructed manually.

\begin{figure*}[H]
\centering
\begin{subfigure}[t]{0.49\textwidth}
\centering
\includegraphics[width=\linewidth]{figures/ef004_20d_all.pdf}
\caption{Discovery instances 1--5 (EF004).}
\label{fig:coco_discovery}
\end{subfigure}
\hfill
\begin{subfigure}[t]{0.49\textwidth}
\centering
\includegraphics[width=\linewidth]{figures/ef005_20d_all.pdf}
\caption{Held-out instances 6--10 (EF005).}
\label{fig:coco_heldout}
\end{subfigure}
\caption{Official COCO target-runtime ECDF views for the same four frozen configurations on all 24 noiseless BBOB functions at \(D=20\). The discovery and held-out panels are produced from disjoint instance sets under the same scientific configuration. The figure is used for condition-level characterization, not for a universal ranking or a post hoc binary failure boundary.}
\label{fig:coco_discovery_heldout}
\end{figure*}

The principal empirical observation is not a single global winner. The relative separation among methods varies with target difficulty and problem class, and the held-out panel preserves the same broad heterogeneity seen during discovery rather than collapsing to a stable total ordering. This is precisely why the paper reports ECDF/target-runtime behavior together with condition metadata instead of reducing the campaign to one average rank.

\subsection{Validated scale of the campaign}

Table~\ref{tab:campaign_scale} summarizes the promoted computational products. EF003 supplies fixed-budget cross-dimensional scaling; EF004 supplies algorithm-separated COCO target-runtime logging on discovery instances; EF005 repeats the target-runtime protocol on disjoint held-out instances.

\begin{table}[H]
\centering
\caption{Validated computational products used in the manuscript.}
\label{tab:campaign_scale}
\footnotesize
\begin{tabular}{p{0.11\linewidth}p{0.21\linewidth}p{0.18\linewidth}p{0.35\linewidth}}
\toprule
Freeze & Instance role & Scale & Publication use \\
\midrule
EF003 & Discovery, 1--5 & 1,152 dimension--algorithm--function--budget cells & Descriptive fixed-budget scaling at 100D, 300D, and 1000D \\
EF004 & Discovery, 1--5 & 14,400 checkpoint rows; 64 algorithm-labelled COCO roots & Official target-runtime/ECDF characterization \\
EF005 & Held-out, 6--10 & 14,400 checkpoint rows; 64 algorithm-labelled COCO roots & Held-out target-runtime/ECDF characterization \\
\bottomrule
\end{tabular}
\end{table}

Across EF004 and EF005, the combined validation logic required exact agreement between the final ObjectiveRecorder and COCO evaluation counts. Joint \texttt{cocopp} postprocessing completed for both campaigns. These checks do not establish algorithmic superiority; they establish that the performance views in Figure~\ref{fig:coco_discovery_heldout} are tied to validated evaluation accounting and algorithm-separated COCO roots.

\subsection{Fixed-budget and target-runtime views}

The study retains two complementary views. Fixed-budget analysis asks what has been achieved after a declared number of evaluations. Target-runtime analysis asks how many evaluations are required to reach a declared target, with unreached targets retained as unsuccessful observations. A method can be comparatively strong under one view and weaker under another without contradiction.

For target-runtime analysis, COCO's logger and postprocessing define the authoritative attainment semantics. Exact target runtimes are not reconstructed from sparse checkpoints, and private optimum values are not inferred from undocumented internals. For fixed-budget scaling, FE/D remains explicit so that the same nominal budget is not confused with equal wall-clock cost.

\subsection{Interpretation boundary}

The held-out campaign was designed to prevent the discovery instances from serving as their own sole confirmation. However, the original failure-frontier protocol required an exact numerical reliability threshold and target subset to be frozen before the held-out results were inspected. Those numerical choices were not fully frozen before EF005 execution. We therefore do not retrofit a binary frontier after observing the confirmation data.

This restriction is scientifically consequential. The experiments support condition-level target-runtime and scaling characterization across the evaluated BBOB conditions; they do not support a universal failure boundary, a universal optimizer ranking, or a mechanism-repair claim derived solely from deterioration. R12 was therefore closed as \texttt{NOT\_WARRANTED}, and no new optimizer is introduced.

\subsection{Cost and fairness}

Equal objective-evaluation budgets improve comparability but do not equalize computational cost. CMA-ES maintains covariance-related state, DE operates on a population, NM maintains simplex geometry, and Random Search has little optimizer-side state. Objective evaluations and optimizer overhead are therefore treated as different resources.

The comparator panel is similarly bounded. PSO and GA were not added merely to enlarge the experiment because their publication implementation/configuration provenance had not passed the same acceptance path when the panel was frozen. The paper favors fewer fully specified comparators over a broader but weaker comparison.

\subsection{Result synthesis}

The computational evidence supports three bounded conclusions. First, the same scientific configuration can be executed on disjoint discovery and held-out BBOB instances with exact objective-call accounting. Second, official COCO target-runtime products reveal substantial condition dependence, so a single aggregate rank would discard information central to interpretation. Third, the held-out campaign can validate a characterization claim while still being insufficient for a stronger threshold-dependent frontier claim. The empirical contribution is therefore not a new optimizer; it is an auditable demonstration of how claim strength changes when the experimental protocol is held fixed and the interpretation rule is enforced.

\section{Validated Gaps, Existing Remedies, and Evidence Obligations}
\label{sec:gaps}

\subsection{A gap is not automatically a discovery}

Optimization papers frequently use the term gap for very different situations: an algorithm has not been tested on a problem class, two studies disagree, a limitation is observed under one condition, or an existing mechanism appears unable to repair a reproduced failure. These situations require different evidential responses. The atlas therefore treats gap strength as an ordered validation problem rather than a binary label.

At the weakest level, a topic may simply be mentioned as underexplored. Stronger evidence can establish an absence of controlled study, a contradiction between comparable studies, or a reproducible condition-dependent limitation. A mechanism-level gap requires an additional step: evidence that the limitation is plausibly connected to a particular computational mechanism rather than to tuning, implementation, budget, or benchmark selection. A repair opportunity is stronger still because it requires reason to believe that an existing remedy is insufficient and that a specific intervention can address the implicated mechanism.

\subsection{Current evidence obligations}

The frozen status matrix identifies several bounded evidence obligations. For Nelder--Mead, the unresolved question concerns dimension and simplex geometry rather than a universal claim that the method fails in high dimension. For BFGS, nonconvex curvature behavior should be separated across globalization and damping regimes. For L-BFGS, documented nonsmooth limitations motivate controlled extension beyond known counterexample classes rather than rediscovery of the existing limitation.

For simulated annealing, the unresolved finite-budget question couples barrier structure, dimension, and cooling policy. For ant-colony optimization, reports of premature pheromone concentration motivate a state-indicator and remedy audit. In multiobjective optimization, NSGA-II motivates diagnostic mapping of crowding behavior as the number of objectives and front geometry change, whereas MOEA/D and NSGA-III motivate controlled studies of decomposition and reference-direction behavior on difficult Pareto-front geometries.

These entries are not ranked by importance. They illustrate how a broad limitation statement can be converted into a condition-indexed experiment that could either strengthen, weaken, or refine the claim.

\subsection{Existing remedies must be tested before novelty is claimed}

A limitation does not justify a new optimizer if an established mechanism already addresses it. Restart and multistart strategies can alter stagnation behavior. Adaptive parameter control can modify fixed-control failure regimes. Population-size schedules can change the exploration budget. Archives and diversity mechanisms can preserve alternative search information. Trust regions can localize model-based search. Constraint-handling rules can replace poorly scaled penalty formulations. Decomposition, reference directions, indicators, and archive truncation can alter multiobjective selection geometry.

The relevant scientific question is therefore not whether a new operator can be invented, but whether the observed failure persists after credible existing remedies are considered. A proposed mechanism should identify which state variable, information source, proposal geometry, memory rule, selection rule, adaptation process, or constraint model it changes and why that change addresses the reproduced limitation.

\subsection{Why no new optimizer is introduced here}

The project completed a frozen discovery campaign and an independent held-out characterization. However, the exact numerical target subset and reliability threshold required for a confirmatory binary frontier were not frozen before the held-out campaign. The project therefore does not elevate observed deterioration into a post hoc exact frontier.

Without a confirmed frontier, an independently supported mechanism explanation, and a comparison against established remedies, a new optimizer or repair operator would be result-driven. R12 was consequently closed as \texttt{NOT\_WARRANTED}. This is not a negative result for the atlas. It demonstrates that the same framework used to identify research opportunities can also reject premature novelty.

\subsection{Minimal standard for a future repair claim}

A future mechanism-repair study can proceed from the atlas if it satisfies five requirements. First, the failure event and threshold must be frozen before confirmation. Second, discovery and confirmation conditions must be disjoint. Third, the implicated mechanism must have evidence independent of the performance transition itself. Fourth, credible existing remedies must be included as comparators or ablations. Fifth, the intervention should be tested as a transferable mechanism where compatible, for example by comparing \(A\) with \(A+M\) across more than one unrelated algorithm rather than introducing only a new named optimizer.

This design would distinguish a genuine mechanism contribution from retuning. It would also permit competing explanations to be tested directly. Only after such evidence should a general mechanism claim be considered.

\subsection{Evidence obligations as a research frontier}

The atlas therefore reframes part of the optimization research frontier as a set of evidence obligations. Some obligations concern missing benchmark domains, some concern unresolved mechanism--condition relations, some concern contradictory studies, and others concern reproducibility or computational cost. Their common feature is that each can be associated with a minimal resolving action.

This view is deliberately more conservative than a list of open problems. An open problem may be intellectually interesting but underspecified empirically. An evidence obligation identifies what observation is missing and what would change the current evidence state. In this sense, the atlas is intended not only to summarize optimization research, but also to expose where additional evidence would be maximally informative.

\section{Researcher Decision Framework and Research Frontier}
\label{sec:decision}

\subsection{Begin with the problem, not the algorithm name}

The atlas does not support a universal optimizer recommendation. A defensible selection begins by describing the decision problem: variable representation, derivative availability, dimensionality, evaluation cost, constraints, uncertainty or noise, temporal change, number of objectives, and available evaluation budget. These conditions determine which information an optimizer can exploit and which comparison is meaningful.

The first decision is therefore an information decision. When reliable derivatives and appropriate smoothness are available, gradient or curvature information should not be discarded without reason. When derivatives are unavailable or unreliable, derivative-free mechanisms become relevant. Expensive objectives motivate sample-efficient or surrogate mechanisms; mixed, discrete, constrained, robust, dynamic, and multiobjective settings require mechanisms whose state and selection rules represent those structures.

\subsection{Choose comparators by scientific role}

Comparator selection should test an explanation rather than assemble a long list of names. A useful panel normally includes a transparent reference baseline, a mechanistically close comparator, an established method appropriate to the regime, and a specialist when the hypothesis concerns a particular problem structure. Competitors should be selected before outcomes are inspected.

Mechanism fingerprints help identify the closest structural comparator. If a proposed contribution changes parameter memory in DE, comparison only with unrelated swarm algorithms would provide weak evidence about the claimed mechanism. Parent-versus-parent-plus-mechanism comparisons, together with established adaptive variants, are more informative. Conversely, a broad application study may require representatives from genuinely different mechanism classes.

\subsection{Match benchmarks to the claim}

Benchmark selection should follow the claimed problem properties. A scalability claim requires a dimension axis; a noise-robustness claim requires controlled noise; a constraint-handling claim requires varying constraint structure; and a mixed-variable claim cannot be established on continuous BBOB alone. Benchmark diversity is useful only when it measures the dimensions of the claim.

The same principle limits aggregation. Continuous, noisy, discrete, dynamic, and multiobjective results should not be averaged into a universal score merely because all are optimization problems. Property-conditioned synthesis is preferable when the benchmark mapping and evidence density support it.

\subsection{Treat budget and cost as separate design variables}

Objective evaluations, wall-clock time, memory, communication, and optimizer-side computation measure different resources. Equal FE/D budgets are appropriate when objective evaluations dominate cost, but they do not make covariance adaptation, population operations, surrogate fitting, and random sampling computationally equivalent. Studies should therefore identify the resource constrained by the application and report secondary costs when they could change the practical conclusion.

For stochastic comparisons, matched instances and multiple seeds remain important. Deterministic algorithms should not acquire artificial replication merely because stochastic competitors use several seeds. The statistical unit should follow the actual source of variation.

\subsection{Use the evidence state to choose the next experiment}

A documented condition may require replication or extension rather than rediscovery. A partial condition calls for an experiment that isolates the unresolved axis. Contradictory evidence calls for a design that reproduces the conditions under which the reports differ. An untested condition calls first for characterization, not a failure claim. Project-validated evidence can support a stronger conditional statement but remains bounded by the frozen protocol.

This creates a practical decision rule: the next experiment should be the smallest experiment capable of changing the evidence state. Large benchmark campaigns are not automatically more informative if they do not resolve the uncertainty that motivated them.

\subsection{Novelty checking at the mechanism level}

A new name is insufficient evidence of algorithmic novelty. A novelty audit should identify the changed state representation, information source, proposal geometry, memory rule, selection rule, adaptation policy, constraint mechanism, or model. The closest prior mechanisms should then be compared directly.

If the contribution is a transferable component, it is often more informative to test \(A\) versus \(A+M\) across multiple compatible parent algorithms than to package \(M\) immediately as an independent optimizer. Such a design tests whether the proposed mechanism, rather than unrelated implementation choices, explains the observed change.

\subsection{Research frontier as unresolved evidence}

The research frontier emerging from the atlas is not a ranked list of fashionable algorithms. It consists of unresolved evidence cells for which the missing observation can be stated precisely. Examples include controlled dimension--geometry mapping for simplex search, globalization-sensitive nonconvex curvature studies, barrier--dimension--cooling interactions in simulated annealing, pheromone-state diagnostics and remedy audits in ant-colony optimization, and geometry-conditioned studies of multiobjective selection mechanisms.

Broader obligations remain in noisy, discrete, mixed-variable, dynamic, robust, learning-assisted, and real-world optimization. Their presence in the census does not imply adequate comparative evidence. These domains should be expanded through separately frozen protocols appropriate to their information structures rather than appended to continuous BBOB results for breadth.

\subsection{Decision principle}

The atlas ultimately supports a simple research principle:
\begin{equation}
\text{condition}\rightarrow\text{mechanism}\rightarrow\text{evidence}\rightarrow\text{experiment}\rightarrow\text{claim}.
\end{equation}
Reversing this order---starting from a desired algorithm or claim and then searching for favorable conditions---weakens both novelty and reproducibility. The framework is intended to keep algorithm choice, experimental design, and claim strength connected to the evidence that can actually support them.

\section{Threats to Validity and Limitations}
\label{sec:limitations}

\subsection{Literature coverage}

No review of optimization can guarantee exhaustive coverage of every named method, application, or publication. The 90-entry census is a structured coverage layer, while the 42 source-backed fingerprints are representative mechanisms rather than a claim of completeness. Relevant work may be missed because of indexing, terminology, language, publication venue, or the rapid appearance of new methods. The atlas therefore treats its scope as auditable and extensible rather than closed.

Publication bias is a related concern. Positive algorithm comparisons are more likely to be emphasized than null or unfavorable results, and benchmark studies can differ substantially in tuning effort. The evidence-state model reduces but cannot eliminate this bias.

\subsection{Taxonomy subjectivity}

Mechanism normalization requires judgment. Fields such as dominant information source, memory, adaptation, diversity control, or failure signal can admit more than one defensible encoding. Structural similarity therefore depends on the selected representation and field weights. Similarity is used as an exploratory aid, not as proof of equivalence, ancestry, or lack of novelty.

The mechanism-loop abstraction is deliberately broad. Its value is comparability, but that breadth can hide domain-specific mathematical structure. Formal convergence properties and specialized problem representations must still be examined in their native theoretical frameworks.

\subsection{Evidence heterogeneity}

Published studies differ in benchmark suites, implementations, parameter tuning, stopping rules, computing environments, and statistical practice. Such heterogeneity limits quantitative pooling. The atlas therefore favors condition-specific synthesis and marks partial or contradictory evidence rather than forcing incomparable studies into a single effect estimate.

The reproducibility registry is also incomplete for some representatives. Pending fields are not interpreted as negative evidence, but they limit conclusions about the reproducibility landscape.

\subsection{Benchmark representativeness}

The project-generated experiment is restricted to noiseless continuous BBOB. BBOB provides a controlled and widely used black-box test environment, but it does not represent discrete, mixed-variable, constrained, dynamic, noisy, expensive, multiobjective, or application-specific optimization in general. Results from the project experiment should not be transferred to those domains.

Even within BBOB, a finite set of dimensions, instances, targets, and budgets samples only part of the possible condition space. Property-group interpretations depend on the validity of the benchmark-property mapping and sufficient observations within each group.

\subsection{Comparator and configuration scope}

The publication comparator panel contains four frozen configurations: Random Search, SciPy differential evolution, bounded SciPy Nelder--Mead, and pycma CMA-ES. This panel spans distinct mechanisms but is not intended to represent the full state of the art. Algorithms without an accepted publication implementation/configuration path were excluded rather than included with weaker provenance.

Performance is configuration-dependent. A result for the frozen DE or CMA-ES configuration is not a result for every variant of that family. The study also does not claim that equal objective-evaluation budgets equalize optimizer-side computational cost.

\subsection{Statistical and frontier limitations}

Discovery and held-out instance sets reduce one source of overfitting, but they do not provide unlimited external validation. More importantly, the exact reliability threshold and target subset required for a binary failure-frontier claim were not frozen before held-out execution. The study therefore does not report such a frontier. This restriction reduces claim strength but protects the distinction between predeclared confirmation and post hoc threshold selection.

No G4/G5 mechanism-repair opportunity was promoted from the frozen evidence. Consequently, the study does not test a new optimizer or repair operator. This limits claims about mechanism causality but avoids a result-driven intervention.

\subsection{Temporal validity}

Optimization research changes rapidly. The source ledger and census represent the evidence available under the review's search and freeze process, not a permanent endpoint. New algorithms, theoretical results, benchmark suites, and independent replications can change individual evidence states. The machine-readable atlas and versioned evidence model are intended to make such updates traceable.

\subsection{Reproducibility versus generalizability}

The frozen workflows strengthen computational reproducibility, but reproducibility and generalizability are distinct. A perfectly reproducible result can remain narrow if the benchmark distribution, implementation, or parameter policy is narrow. Conversely, a broad literature pattern may be difficult to reproduce exactly when original artifacts are unavailable.

The conclusions therefore remain bounded by both dimensions: what can be rerun and what the evaluated conditions can reasonably represent. This boundary informs the conclusions that follow.

\section{Conclusion}
\label{sec:conclusion}

This paper developed an optimization atlas in which algorithm names, computational mechanisms, evidence states, reproducibility, and research gaps are distinct but connected objects. The separation matters because optimization contains both genuine methodological diversity and structural overlap, while empirical conclusions remain sensitive to problem conditions, budgets, implementations, and parameter policies.

The mechanism-first representation organizes algorithms by the state they preserve, information they use, proposal geometry, and the memory, selection, adaptation, diversity, constraint, and model mechanisms that modify search. Within the declared scope, a 90-entry census provides the identity layer, while 42 source-backed fingerprints linked to 46 verified source records provide the representative mechanism layer. The counts serve different purposes and are not interpreted as 90 distinct mechanisms.

The evidence atlas separates documented, partial, contradictory, untested, and project-validated conditions. This prevents absence of study from being labelled as failure and prevents structural similarity from transferring empirical conclusions between algorithms. Theory is retained with its assumptions and guaranteed object rather than collapsed into a generic notion of convergence.

The project-generated BBOB study demonstrates the executable side of the framework. The frozen cross-dimensional campaign provides a validated condition-level characterization across four comparator mechanisms, 24 noiseless BBOB functions, dimensions 5--40, and normalized evaluation budgets. Algorithm-separated COCO logging, independent objective-call accounting, and held-out instances provide an auditable path from execution to interpretation. The experiment does not establish a universal ranking.

The evidence audit also constrains what the paper does not claim. Because the exact numerical decision rule required for a confirmatory binary failure frontier was not fully frozen before held-out execution, no post hoc exact frontier is reported. Because the resulting evidence does not establish a G4/G5 mechanism-level repair opportunity, no new optimizer or repair operator is introduced. Claim strength is limited by the strongest evidence specified and validated before interpretation.

The practical outcome is a decision sequence from condition to mechanism, evidence, resolving experiment, and claim. Future versions can add algorithms, benchmark domains, theoretical results, replications, or application evidence without changing the meaning of existing records. Unresolved cells can be converted into explicit evidence obligations whose minimal resolving experiment is known.

Optimization research does not require every unresolved condition to become a new algorithm. In many cases the more informative contribution is to determine what is known, what remains untested, which mechanism is implicated, and what evidence would distinguish a genuine limitation from an artifact of experimental design. The atlas provides a reproducible structure for making that distinction.

% Subsequent sections are drafted in frozen-evidence order.
\bibliographystyle{IEEEtran}
\bibliography{references}
\end{document}
